\input{./._relpath-to-latexroot.ltx}
\documentclass[\latexroot/\projectname]{subfiles}
\input{\latexroot/@resources/subfile-setup.ltx}

\begin{document}

\section{Introduction}
\whenintegrated{\label{sec:intro}} % integrated doc: SST for labels
\setcounter{page}{0}\pagenumbering{arabic}

Government debt has risen substantially across advanced economies in recent
decades. Yet countries with similar levels of public debt often differ sharply
in who holds that debt. In some countries, government debt is absorbed mainly
by domestic non-bank investors such as pension funds, insurance companies, and
households. In others, a larger share is held by foreign investors or by the
domestic banking system, including private banks and the central bank. This
raises a central question: when a government has more debt to place, what
determines who ends up holding it, and why does this allocation differ across
countries?

This question matters for both fiscal and financial stability. The creditor
composition of public debt affects the strength of the sovereign--bank nexus,
the extent of fiscal space, and the economy's exposure to global investor
sentiment. A government that relies heavily on foreign investors may be more
vulnerable to shifts in global risk appetite, while a government that places a
large share of debt with domestic banks may deepen the connection between
sovereign risk and financial-sector fragility. Similarly, a country with a
large domestic non-bank investor base may be able to absorb public debt with
less pressure on the banking system or foreign funding.

Despite its importance, the allocation of sovereign debt across creditor
sectors remains insufficiently understood. Existing work often studies one
creditor sector at a time, or documents differences in debt-holder composition
without providing a joint account of why the split across sectors takes the
form it does. What is missing is a framework that explains how a given stock of
government debt is allocated across domestic non-banks, foreign investors, and
the banking sector, and how that allocation depends on both the level of debt
and the country's domestic financial structure.

This paper addresses that gap by developing a cost-minimizing model of
sovereign debt allocation. The government must place a given amount of debt and
chooses how much to allocate to three sectors: domestic non-banks, foreign
investors, and the banking sector. Each sector differs in its effective
capacity and marginal placement cost. Domestic non-bank capacity varies across
countries with the depth of the local financial system. Foreign demand is
limited by a common exposure ceiling. The banking sector has no hard ceiling,
but greater reliance on banks carries increasing quasi-fiscal and
financial-stability costs. The model therefore links cross-country differences
in debt-holder composition to differences in domestic financial structure and
predicts how marginal debt absorption changes as total debt rises.

We structurally estimate the model using panel data on debt composition across
countries and over time. The estimation identifies distinct borrowing limits
across creditor types and provides a tractable framework to interpret
variation in who holds government debt.

% Smart bibliography: Only include bibliography if standalone AND has citations
\smartbib

\end{document}
